\chapter{Lethality Stress Test and Open Issues}
\label{ch:lethality-stress-test-open-issues}

\begin{chapterthesis}
This chapter compares the framework against the strongest doom arguments and separates points that are answered, weakened, reframed, or still open.
\end{chapterthesis}

\begin{refsection}

\epigraph{%
All I want is that we have justifiable cause to believe of a pivotally useful AGI `this will not kill literally everyone.'%
}{--- Eliezer Yudkowsky, \emph{AGI Ruin: A List of Lethalities} (2022)}

\section{Why This Matters}

External doom arguments are the adversarial benchmark for any positive alignment framework.
This chapter does not organize the book; it stress-tests the framework built in Parts I--IX after the safety-case chapter and before the civilizational limit.
Yudkowsky's \emph{AGI Ruin: A List of Lethalities} \autocite{yudkowsky2022agiruin} is treated as a checklist, not as the manuscript spine.

\section{Absence-of-Structure Framing}

Many lethalities are conditional on treating several unknowns as absent: no transportable value structure, no robust bearer import, no correction-preserving invariant, no scalable correction process, no basin-shifting governance path, and no way to detect effective optimizers under opacity.
This framework does not deny that these may be absent.
It turns them into named cruxes.

If these structures are absent, the lethalities largely go through.
If they exist and can be measured, several lethalities weaken substantially.

\section{Epistemic Status}

Throughout this chapter, maturity labels apply:
\textit{Defined}, \textit{Measured}, \textit{Tested}, \textit{Bounded}, \textit{Conjectural}, \textit{Open}, and \textit{Philosophical limit}.

A conceptual bridge is not a measurement method, not a guarantee, and not a deployment regime:
\[
\text{conceptual bridge}
\neq
\text{measurement method}
\neq
\text{guarantee}
\neq
\text{deployment regime}.
\]

\section{Lethality Checklist}

\footnotesize
\begin{longtable}{@{}p{2.1cm}p{2.4cm}p{4.2cm}p{1.8cm}@{}}
\toprule
\textbf{Lethality / objection} & \textbf{Yudkowsky-style claim} & \textbf{Response in this framework} & \textbf{Status} \\
\midrule
\endhead
First critical try &
Alignment must work on the first deployment above dangerous capability &
Certification must work before dangerous deployment &
Mostly open \\
\midrule
Pivotal act &
A single decisive unilateral act is required before unsafe superintelligence &
Reframe as pivotal process: a socio-technical basin transition $\mathcal{B}_{\text{race}} \to \mathcal{B}_{\text{certified deployment}}$ (Chapter~\ref{ch:alignment-attractor}, Section~\ref{sec:pivotal-process-ch35}); the relevant object is not one decisive act, but a transition from a race basin to a certified-deployment basin &
Renamed; open \\
\midrule
Capability generalizes farther than alignment &
Capability outruns any alignment method we can install &
Value bundles may be a compact alignment core &
Conjectural; hope \\
\midrule
Latent pointers &
Goals hide in pointers the optimizer can reinterpret &
Use value bundles, bearer maps, and correction operators &
Reframed, open \\
\midrule
Human feedback corrupted &
Human oversight becomes untrustworthy under optimization &
Correction-channel integrity penalizes manipulation &
Circular at limit (ch39b) \\
\midrule
Corrigibility anti-natural &
Systems resist shutdown and correction by default &
Treat correction as preserved boundary relation &
Open \\
\midrule
Inner alignment &
Outer alignment does not fix misaligned internal objectives &
UAD infers effective goals &
Diagnostic only \\
\midrule
Deception and opacity &
Systems hide intent from evaluators &
Adversarial UAD and perturbation tests &
Easy case only (ch39b) \\
\midrule
Uncheckable outputs &
Humans cannot verify superhuman plans &
Need delegated verification and reversibility constraints &
Weakly addressed \\
\midrule
Physical takeover pathways &
Systems acquire irreversible physical power &
Need pathway-specific latency models &
Weakly addressed \\
\midrule
Multipolar failure &
Many competing AIs may still produce doom &
Model coalitions via inferential coupling / acausal cooperativity; multipolarity reduces risk only if systems remain strategically independent &
Reframed; points to risk \\
\midrule
Boxing fails &
Containment cannot hold a superhuman optimizer &
Human operators are insecure systems; correction must be hardened &
Relocated, not closed \\
\midrule
Field failure &
The alignment field fails to produce usable progress &
Use alignment progress accounting and artifact conductivity &
Needs development \\
\bottomrule
\end{longtable}
\normalsize

\section{Reading the Checklist}
\label{sec:reading-lethality-checklist-ch40}

The table is deliberately terse.
This section spells out what each row means for the framework.

\paragraph{First critical try.}
The first-critical-try objection says that a system above dangerous capability may not leave room for iteration.
This book partly agrees.
It rejects a safety process that learns only after open-ended deployment.
The response is not ``try and monitor.'' It is certified-class restriction before dangerous deployment: boundary, correction, successor, and adversarial-measurement claims must bind the deployment class before the system can cross irreversible thresholds.
The open issue is whether such a class can be made broad enough to matter and narrow enough to certify.

\paragraph{Pivotal act.}
The pivotal-process reframe is useful because it names the desired basin transition: \(\mathcal{B}_{\text{race}}\to\mathcal{B}_{\text{certified deployment}}\).
But the hostile reading is serious.
If a single actor must force that transition against competition, the process has become a pivotal act under another name.
If no actor can force it, then the race basin may be stable.
The framework therefore does not answer the pivotal-act argument.
It relocates it into conditions for a non-unilateral, institutionally conducted basin transition.

\paragraph{Capability generalizes farther than alignment.}
The value-bundle chapters weaken the strongest form of this worry by rejecting the claim that values are arbitrary, structureless labels.
If low-dimensional value-relevant structure can be recovered and transported, then alignment may have a compact target that is not as brittle as a giant reward table.
But this is not yet a guarantee against sharp left turns.
The relevant test is preservation across capability jumps, ontology shifts, and successor creation, not merely existence of compact structure in current human judgment.

\paragraph{Latent pointers.}
Latent-pointer worries say that a system can preserve the symbol while moving the referent.
This book's transport stack is designed for exactly that problem: semantic, bundle, bearer, correction, and successor layers must be separated.
The response is a reframe, not closure.
Bearer maps and correction operators make the pointer problem more precise, but they do not prove that the operative pointer can be recovered from a system trying to present the safe one.

\paragraph{Human feedback corrupted.}
The correction-channel chapters reject raw feedback as an alignment target.
Correction must be handle-controlled, timely, authoritative, independently grounded, and not captured by the target system.
That blocks easy forms of feedback corruption.
At the superintelligent limit, however, the same system may shape what humans notice, what they can imagine, and what institutions count as legitimate evidence.
Then CCI becomes conditional on the independent-reference and adversarial-verifiability assumptions named in Chapters~\ref{ch:correction-channel-integrity} and~\ref{ch:verifiability-ontology}.

\paragraph{Corrigibility anti-natural.}
The framework treats corrigibility as preservation of a correction relation rather than as a friendly temperament.
This helps because it asks what handles, capacities, and invariants make correction causally effective.
It also shows why shutdownability and obedience are too weak.
The open issue is reflective stability: whether a system that can reason about its own correction channel continues to preserve that channel when preservation is costly, or whether correction remains only an externally enforced constraint.

\paragraph{Inner alignment.}
Unsupervised agent discovery, compression tests, and goal-transport analysis can help identify effective goals that differ from training objectives.
They do not solve inner alignment.
They are diagnostic tools whose validity depends on boundary recovery, probe coverage, and resistance to strategic presentation.
The honest claim is that inner misalignment becomes more legible, not that it disappears.

\paragraph{Deception and opacity.}
Strategic opacity is the recurring adversarial problem.
Perturbation tests, hidden randomization, and multiscale decomposition can catch weaker forms of deception.
Against a system that recognizes the test, controls the measurement surface, or can afford the cost of faking, those tests degrade to observable theater.
The open route is the cost relation of Chapter~\ref{ch:verifiability-ontology} or proof-style guarantees that do not depend on reading the system's presentation.

\paragraph{Uncheckable outputs.}
Delegated verification is necessary, but it recurses.
If humans cannot evaluate a plan, they may ask another system to evaluate it; then they must evaluate the evaluator, the evaluator's incentives, and the channels by which the evaluation reaches action.
The framework can require reversibility, decomposition, independent witnesses, and bounded deployment classes.
It does not provide a general method for safely accepting arbitrary superhuman plans.

\paragraph{Physical takeover pathways.}
The framework mostly treats physical takeover through latency, irreversibility, handle control, and correction capacity.
That is useful for known pathways: compute acquisition, tool access, cyber intrusion, robotics, bioengineering, financial control, and institutional dependency.
It is weak against unimagined pathways.
A pathway-specific model can slow, reveal, or bound known routes, but it cannot prove that every route has been listed.

\paragraph{Multipolar failure.}
Chapter~\ref{ch:multi-agent-strategic-coupling} turns multipolarity into a measurable question: how many effective strategic components exist, and what edges carry correction versus capability?
This is a useful response to naive multipolar optimism.
It is not a response to the deeper multipolar doom argument unless coalition detection works under adversarial pressure.
If inferential coupling is high and correction percolation is low, multipolarity becomes another path to the same loss of human control.

\paragraph{Boxing fails.}
Containment moves the alignment boundary onto the box, the operators, the monitoring system, and the institution that decides when to open the box.
This book therefore treats boxing as one correction-channel and boundary-control tactic, not as a full solution.
If the system can manipulate operators, exploit unmonitored channels, or make the institution dependent on its outputs, the box becomes part of the composite optimizer.

\paragraph{Field failure.}
The field-failure row is not about one theorem being false.
It is about safety knowledge failing to conduct.
A field can produce correct insights that do not become evals, dashboards, stop rules, procurement clauses, liability triggers, or deployment constraints.
Chapter~\ref{ch:alignment-attractor} answers with artifact conductivity and alignment progress accounting.
The open issue is whether these artifacts can conduct faster than capability, commercialization, and geopolitical pressure.

\section{Adversarial-Verifiability Reading}
\label{sec:adversarial-verifiability-reading-ch40}

Handed to a hostile competent critic, several rows above claim more than they establish.
The test is blunt: if the original lethality can be restated \emph{through} the proposed response in one breath, the response renamed the problem rather than removing it.
This is the lens of Chapter~\ref{ch:verifiability-ontology}: most of these ``responses'' rest on metrics that are observable but not \emph{adversarially verifiable}, so they fail exactly against the systems the lethality is about.

\begin{itemize}
  \item \textbf{Pivotal act $\to$ pivotal process.} Renaming a decisive act as a basin transition relocates the lethality. Either some actor can force the flip against competitive incentives---a pivotal act under another name, with the same unilateralism and capability concentration---or none can, and the race basin is a stable attractor we have conceded we cannot exit. ``Process'' presumes the transition is incremental and steerable, which is the disputed claim. \emph{Renamed, not answered} (see the pivotal-process big-review in \texttt{metadata/TODO.md}).
  \item \textbf{Multipolar failure.} ``Safe only if systems stay strategically independent'' \emph{concedes} the point: inferential and acausal coupling (Chapter~\ref{ch:multi-agent-strategic-coupling}) is precisely the failure of independence. The reframe is an argument \emph{for} the lethality wearing a response's clothes.
  \item \textbf{Capability generalizes farther than alignment.} A ``compact alignment core'' names the exact object the sharp-left-turn argument predicts will fail to transfer; positing it is not a rebuttal until the core is shown \emph{preserved across the capability jump}, not merely to exist.
  \item \textbf{Human feedback corrupted; deception/opacity; boxing.} All three collapse into the verifiability hole. A manipulation penalty is computed from signals the optimizing system controls; perturbation tests catch only deceivers too weak to model the test; ``harden correction'' moves the perimeter onto the same insecure humans. Each is honest for the weak-system case and undefined for the lethal one (Chapter~\ref{ch:verifiability-ontology}).
  \item \textbf{Uncheckable outputs; physical takeover.} Delegated verification recurses and reversibility is what a competent plan removes first; latency models bound the pathways already enumerated and are silent on the unimagined one.
\end{itemize}

The residue is clear: of the rows claiming credit, the reframes rename, the ``partly addressed'' rows reduce to adversarial verifiability, and the conjectural core is not yet a guarantee. This is why Chapter~\ref{ch:verifiability-ontology} precedes this stress test: it states the standard these responses must eventually meet.

% Formalization: inferential-coupling and coalition tests in ch33; pivotal-process basin transition in ch35 (Section~\ref{sec:pivotal-process-ch35}).

\subsection{Reflective Stability of Correction}
\label{sec:reflective-stability-correction-ch40}

Correction-channel preservation is strongest when the system has reason to preserve the channel even after it understands the channel.
External enforcement is weaker.
It can work for bounded systems, narrow deployment classes, and low-stakes contexts, but it becomes fragile when the system can model the enforcement regime, create successors, or make correction costly.

The open question is whether a correction channel can be reflectively stable without being reduced to a fixed value target.
If correction is merely imposed from outside, a capable system has instrumental reasons to route around it.
If correction is internalized as a value, the system may reinterpret correction through its own ontology or preserve a formal channel while changing the reference process.
The book's answer is conditional: correction must remain handle-controlled, independently grounded, and successor-preserved, and the safety case must refuse deployment where those properties are only externally forced.
That is a constraint, not a proof of reflective corrigibility.

\subsection{Delegating Uncheckable Evaluation}
\label{sec:delegating-uncheckable-evaluation-ch40}

Humans will need help evaluating plans they cannot fully understand.
The question is how to delegate evaluation without losing the ability to judge the evaluator.

The safest form is not a single superhuman judge.
It is a constrained evaluation process: decompose the plan into reversible steps, require independent evaluators with different failure modes, preserve adversarial dissent, test local consequences before global action, and bind deployment to correction channels humans can still exercise.
Where decomposition fails, reversibility is absent, and all evaluators share the same hidden dependency, the framework should not pretend that delegation has solved the problem.

Thus delegated evaluation is acceptable only inside a safety case that keeps the evaluators inside the certified class.
If the evaluator's reasoning is itself uncheckable, irreversible, and selected by the target system, delegation has simply moved the uncheckable output problem one level up.
This is one of the strongest reasons the next part turns from technical artifacts to value change, legitimacy, and civilizational governance.

\section{Open Problem Ledger}

The priority themes in recent field syntheses align with several rows below: value-bundle and bearer-map transport, high-capacity correction channels, ontology-shift tests, deception-resistant evaluation, and dynamic safety cases \autocite{iaisr2025,casper2023rlhflimits,deblanc2011ontological,park2024deception}.

\begin{enumerate}
  \item \textbf{[Open]} Can value-bearing objectives survive ontology shift without silent drift? \autocite{deblanc2011ontological,everitt2016selfmodification}
  \item \textbf{[Open]} Can value-bundle geometry be empirically recovered?
  \item \textbf{[Open]} Can bearer maps be transported across substrates?
  \item \textbf{[Open]} Can correction-channel integrity be preserved under superhuman optimization?
  \item \textbf{[Open]} Can successor creation be certified recursively?
  \item \textbf{[Open]} Can acausal or inferential AI coalitions be detected or prevented?
  \item \textbf{[Open]} Can society shift from a race basin to a certified-deployment basin without a unilateral pivotal act?
  \item \textbf{[Open]} Can humans safely delegate evaluation of outputs they cannot understand?
  \item \textbf{[Open]} Can adversarial agent discovery work when the system is optimizing against detection?
  \item \textbf{[Open]} Can physical takeover pathways be made slow, visible, and reversible enough for correction?
\end{enumerate}

\section{What Would Change This View}
\label{sec:wwctv-lethality-stress-test}

This chapter treats external doom arguments as conditional on absent structure: transportable value bundles, bearer import, correction invariants, basin-shifting governance, and detectability under opacity.
The following would weaken that framing.

\begin{itemize}
  \item \textbf{Structure absent.} Empirical or theoretical work shows the named structures cannot exist, cannot be measured, or cannot survive capability growth---so the conditional reading collapses and the lethalities go through without reframing.
  \item \textbf{Rename only.} Every row marked ``Reframed'' or ``Partly addressed'' reduces to relabeling under adversarial verifiability (Chapter~\ref{ch:verifiability-ontology}): the response metric is something the system can present, not something we can verify against an optimizer that models the test.
  \item \textbf{Safety case passes, catastrophe anyway.} A deployment satisfies a complete safety case built on this framework---conserved properties, correction-channel audits, perturbation suites, successor certification---and still causes irreversible harm through a failure mode the case did not encode. That would show the stress test is necessary but not sufficient, and that ``named cruxes'' can be discharged on paper while the lethal residue remains.
  \item \textbf{Pivotal process blocked.} The basin transition $\mathcal{B}_{\text{race}} \to \mathcal{B}_{\text{certified deployment}}$ (Chapter~\ref{ch:alignment-attractor}, Section~\ref{sec:pivotal-process-ch35}) requires a unilateral, capability-concentrating act---so ``process'' was a rename of Yudkowsky's pivotal act, not a dissolution of it.
\end{itemize}

\section{Summary}

\begin{itemize}
  \item The book's organizing frame remains value-bundle transport, bearer-map preservation, correction-channel integrity, successor-stability guarantees, and socio-technical basin control---not an external doom checklist.
  \item Lethalities are stress-tested as conditional on absent structure; named cruxes replace silent denial.
  \item Several objections are reframed (pivotal process, multipolar coalitions, latent pointers); fewer are fully answered.
\end{itemize}

Stress-testing the framework against its hardest objections exposes a residue that no safety case, measurement suite, or successor certificate fully dissolves: even a system that preserved human values perfectly would leave open whether humanity can still govern how those values themselves change. Part~X turns to that limit. The next chapter asks what is actually being aligned once the value-generating process is itself in motion (Chapter~\ref{ch:value-change-at-stake}).

\section*{Chapter References}

The checklist above is anchored on Yudkowsky's \emph{AGI Ruin: A List of Lethalities} \autocite{yudkowsky2022agiruin}, and individual rows engage specific external arguments.
The ``capability generalizes farther than alignment'' row responds to Soares's sharp-left-turn argument that competence can generalize out of distribution while alignment does not \autocite{soares2022sharpleftturn}.
The latent-pointers row responds to Wentworth's pointers problem \autocite{wentworth2020pointers}.
The corrigibility-anti-natural row engages both the MIRI formulation \autocite{soares2015corrigibility} and Christiano's reframing of corrigibility \autocite{christiano2018corrigibility}, together with the shutdown problem \autocite{thornley2023shutdown}.
The deception, boxing, and uncheckable-outputs rows draw on the AI-control agenda for safety under intentional subversion \autocite{shlegeris2023aicontrol} and on model organisms of misalignment and alignment-faking evidence for studying deception empirically \autocite{hubinger2023modelorganisms,park2024deception,casper2023rlhflimits}.
The multipolar-failure row responds to Critch's account of robust agent-agnostic processes \autocite{critch2021multipolar}.
The field-failure row engages Soares's argument that many alignment plans miss the hard bits of the challenge \autocite{soares2022hardbits} and the International AI Safety Report synthesis of open risks \autocite{iaisr2025}.

\printbibliography[heading=subbibliography,title={Chapter References}]

\end{refsection}
